% ===== SECTION 2: COMPANY OVERVIEW =====
\section{\texorpdfstring{\color{MidnightBlue}Company Overview}{Company Overview}}

\subsection{Organization Summary}

\begin{tabularx}{\textwidth}{|>{\bfseries}l|X|}
\hline
\rowcolor{tableHeader}\multicolumn{2}{|l|}{\textcolor{white}{\textbf{Organization Summary}}} \\
\hline
Legal Entity Name          & Fejost LLC dba Sentry Management Solutions \\
\hline
Year Established           & 2016 \\
\hline
Type of Ownership          & Limited Liability Company (LLC) \\
\hline
Parent Company             & None --- Sentry is an independent, privately held company \\
\hline
Headquarters               & \placeholder{Address} \\
\hline
Total Full-Time Employees  & \placeholder{Number} \\
\hline
Primary Contact            & Max Zinner, \placeholder{Title} \\
\hline
Phone / Email              & \placeholder{Phone} / \placeholder{Email} \\
\hline
Website                    & www.sentryms.com \\
\hline
\end{tabularx}

\vspace{0.5cm}

\subsection{Sentry's Role in Student Transportation}

\placeholder{2--3 paragraphs describing Sentry's positioning as a technology and system management provider for student transportation. Emphasize that Sentry is NOT a bus operator but rather provides the technology platform and management services that bus operators and state agencies rely on for routing, dispatch, fleet management, video monitoring, and analytics.}

\begin{sidebar}{Core Capabilities}
\begin{description}[style=nextline, leftmargin=0.2cm]
    \item[\faIcon{route} Routing Optimization:] \placeholder{Describe routing engine, algorithm capabilities, multi-zone optimization}
    \item[\faIcon{headset} Real-Time Dispatch:] \placeholder{Describe dispatch platform, GPS tracking, ETAs}
    \item[\faIcon{video} Video Violation Detection:] \placeholder{Describe camera systems, AI detection, violation processing pipeline}
    \item[\faIcon{chart-line} Analytics \& Reporting:] \placeholder{Describe dashboards, KPI tracking, cost analysis tools}
    \item[\faIcon{bus} Fleet Management:] \placeholder{Describe preventive maintenance tracking, vehicle lifecycle management}
    \item[\faIcon{mobile-alt} Parent \& Driver Communication:] \placeholder{Describe parent app, driver app, real-time notifications}
\end{description}
\end{sidebar}

\subsection{Relevant Experience}

\placeholder{2--3 paragraphs describing Sentry's most relevant engagements. Reference MTA paratransit scale (10,000+ daily trips), school district work, and any video monitoring deployments. Emphasize experience with multi-vendor, multi-zone environments analogous to RIDE's STS structure.}

\begin{tabularx}{\textwidth}{|>{\bfseries}l|X|}
\hline
\rowcolor{tableHeader}\multicolumn{2}{|l|}{\textcolor{white}{\textbf{\placeholder{Key Engagement Title}}}} \\
\hline
Client              & \placeholder{Client Name} \\
\hline
Contract Value      & \placeholder{\$X} \\
\hline
Fleet / Scale       & \placeholder{Description} \\
\hline
Daily Volume        & \placeholder{Description} \\
\hline
Technology Deployed & \placeholder{Description of systems provided} \\
\hline
Service Area        & \placeholder{Geography} \\
\hline
\end{tabularx}

\subsection{Understanding of Rhode Island's STS}

Sentry has carefully studied the current structure of RIDE's Statewide Transportation System, established by the General Assembly in FY~2009 under RIGL~\S~16-21.1 and phased in over three years (14~LEAs in 2010, 26 in 2011, all LEAs by 2012). The system today operates across seven zones serving approximately 3,812 students through 334 routes coordinated by the current Statewide System Manager, \textbf{Transpar}, on behalf of 49~LEAs plus out-of-state placements. RIDE's own \textit{Statewide Transportation System Overview} (October~2024) demonstrates that this system has grown by \textbf{34\%} in ridership since 2015---from 2,850 to 3,812 students---even as total Rhode Island student enrollment has \textit{declined} by 4\% (141,959 to 136,154). The system's FY~2025 cost exceeds \textbf{\$41~million}.

\vspace{0.3cm}

This demand growth is driven by three distinct student populations, each with unique routing and compliance requirements:

\begin{itemize}
    \item \textbf{PCCT (Private, Charter, Career \& Technical)} --- 50.87\% of riders. Regional boundaries under RIGL~\S~16-21.1-2 apply to this population. The Meeting Street example illustrates the consolidation advantage: statewide routing reduced buses serving that single school from 28 to 13---a 54\% reduction.
    \item \textbf{Special Education} --- 24.29\% of riders (234 of 334 routes). Required by federal law (IDEA); regional boundaries do \textit{not} apply. Routes must accommodate wheelchair accessibility, 1:1 aides, and IEP-mandated maximum ride times.
    \item \textbf{Homeless \& Foster Care (McKinney-Vento)} --- 24.84\% of riders. Required by federal law; regional boundaries do \textit{not} apply. Addresses change frequently, requiring dynamic route adjustment.
\end{itemize}

\vspace{0.3cm}

We are particularly attentive to the three bills passed into law in 2025 as a result of the Student Transportation Workgroup's recommendations:

\begin{itemize}
    \item \textbf{Increased Van Size} --- The expanded 10-passenger definition presents an opportunity to convert qualifying routes from higher-cost school buses to vans, reducing per-route costs while maintaining service quality. RIDE has already begun ``rightsizing'' by increasing vans and minibuses; 147 of 334 vehicles are already vans.
    \item \textbf{Temporary School Bus Driver Licenses} --- Abbreviated CDL training for border-state residents can accelerate recruitment pipelines, addressing the driver shortage that competes with Uber, DoorDash, RIPTA, FedEx, and Amazon for the same labor pool. Notably, RI's unique ``white card'' requirement---a multi-day course at CCRI plus DMV driving tests---remains a barrier for out-of-state CDL holders even with the new temporary license provision.
    \item \textbf{Live Digital Video Violation Detection Systems} --- The July 1, 2027 mandate for new buses (and July 1, 2032 for all buses) creates an immediate need for a technology partner capable of deploying, managing, and processing video violation data at scale.
\end{itemize}
